% ---------------------------------------------------------------------
%  CHƯƠNG 2 — CƠ SỞ LÝ THUYẾT VÀ CÔNG NGHỆ
% ---------------------------------------------------------------------
\chapter{Cơ sở lý thuyết và công nghệ}
\label{ch:coso}

\begin{muctieu}
Chương này trình bày các kiến thức nền tảng gồm thuật toán và mô hình trí
tuệ nhân tạo được dùng, cùng danh mục công cụ và thư viện để xây dựng hệ thống.
\end{muctieu}

\section{Kiến thức nền tảng}

\begin{chuy}
\goiy{trình bày ngắn gọn lý thuyết của thuật toán hay mô hình mà em dùng,
ví dụ phân lớp, mạng nơ-ron, cây quyết định, thuật toán tìm kiếm heuristic.
Chỉ nêu phần liên quan trực tiếp tới đồ án, kèm công thức nếu cần. Đây không
phải chương chép lại giáo trình.}
\end{chuy}

\begin{dinhnghia}[Phân lớp văn bản]
Phân lớp văn bản \en{(text classification)} là bài toán gán cho một đoạn văn
bản $x$ một nhãn $y$ trong tập nhãn hữu hạn $\mathcal{Y}$. Bộ phân lớp học
một hàm $f: x \mapsto y$ từ dữ liệu huấn luyện đã gán nhãn.
\end{dinhnghia}

\begin{mau}
Một hướng tiếp cận phổ biến là biểu diễn văn bản theo mô hình túi từ
\en{(bag of words)} kết hợp trọng số TF-IDF, sau đó dùng bộ phân lớp tuyến
tính. Trọng số TF-IDF của từ $t$ trong văn bản $d$ được tính theo công thức
\begin{equation}
\mathrm{tfidf}(t, d) = \mathrm{tf}(t, d) \cdot \log\frac{N}{1 + \mathrm{df}(t)},
\label{eq:tfidf}
\end{equation}
trong đó $\mathrm{tf}(t,d)$ là tần suất từ $t$ trong văn bản $d$, $N$ là
tổng số văn bản và $\mathrm{df}(t)$ là số văn bản có chứa từ $t$. Thừa số
logarit trong~\eqref{eq:tfidf} làm giảm trọng số của những từ xuất hiện ở
hầu hết văn bản, vốn ít giá trị phân biệt.
\end{mau}

\begin{ghinho}
Với dữ liệu vừa và nhỏ, các mô hình tuyến tính như hồi quy logistic hay máy
vec-tơ hỗ trợ thường cho kết quả tốt mà chi phí huấn luyện thấp, rất hợp với
phạm vi một đồ án học phần.
\end{ghinho}

\section{Công cụ và thư viện}

\begin{chuy}
\goiy{liệt kê ngôn ngữ, framework và thư viện chính, ghi rõ phiên bản và vai
trò từng thành phần. Bảng phiên bản chính là thứ giúp người khác chạy lại
được sản phẩm của em.}
\end{chuy}

\begin{table}[H]
\centering
\caption{Công cụ và thư viện sử dụng trong đồ án.}
\label{tab:congcu}
\begin{tabular*}{\textwidth}{@{\extracolsep{\fill}}>{\bfseries\color{daunavy}}l >{\color{daumaroon}}l l}
\toprule
\rowcolor{gray100}\bfseries\color{daunavy}Thành phần & \bfseries\color{daumaroon}Phiên bản & \bfseries\color{daunavy}Vai trò \\
\midrule
Python       & 3.11 & Ngôn ngữ lập trình chính \\
scikit-learn & 1.4  & Huấn luyện và đánh giá mô hình học máy \\
pandas       & 2.2  & Đọc và xử lý dữ liệu dạng bảng \\
underthesea  & 6.x  & Tách từ tiếng Việt \\
Flask        & 3.0  & Xây dựng giao diện web demo \\
\bottomrule
\end{tabular*}
\end{table}

\section{Các công trình liên quan}

\begin{chuy}
\goiy{tóm tắt hai đến ba cách làm sẵn có cho bài toán của em, kèm trích dẫn
bằng lệnh \code{cite} với khoá khai báo trong file \code{refs.bib}. Nêu rõ em
kế thừa gì và làm khác gì. Mục này không bắt buộc nhưng làm báo cáo có chiều
sâu hơn hẳn.}
\end{chuy}

\begin{mau}
Các nghiên cứu về phân tích cảm xúc tiếng Việt thường đi theo hai hướng.
Hướng thứ nhất dùng đặc trưng thủ công kết hợp mô hình tuyến tính, ưu điểm
là nhẹ và dễ giải thích~\cite{HoangKiem}. Hướng thứ hai dùng mô hình học sâu
tiền huấn luyện cho độ chính xác cao hơn nhưng đòi hỏi tài nguyên tính toán
lớn~\cite{Krizhevsky2012}. Đồ án này
chọn hướng thứ nhất làm cơ sở, đồng thời so sánh với một mô hình mạnh hơn để
thấy rõ đánh đổi giữa độ chính xác và chi phí.
\end{mau}
